\documentclass[11pt]{amsart}
\usepackage{amsmath,amssymb,amsthm}
\usepackage[margin=1.05in]{geometry}
\raggedbottom

\numberwithin{equation}{section}
\newtheorem{theorem}{Theorem}[section]
\newtheorem{proposition}[theorem]{Proposition}
\newtheorem{lemma}[theorem]{Lemma}
\newtheorem{corollary}[theorem]{Corollary}
\theoremstyle{definition}
\newtheorem{definition}[theorem]{Definition}
\theoremstyle{remark}
\newtheorem{remark}[theorem]{Remark}
\newtheorem{example}[theorem]{Example}

\newcommand{\Z}{\mathbb{Z}}
\newcommand{\Q}{\mathbb{Q}}
\newcommand{\R}{\mathbb{R}}
\newcommand{\N}{\mathbb{N}}
\newcommand{\C}{\mathbb{C}}
\newcommand{\F}{\mathbb{F}}
\newcommand{\T}{\mathbb{T}}
\newcommand{\hZ}{\widehat{\Z}}
\newcommand{\cA}{\mathcal{A}}
\newcommand{\cO}{\mathcal{O}}
\newcommand{\cP}{\mathcal{P}}
\newcommand{\Idl}{\mathbb{I}}
\newcommand{\Ns}{\mathrm{N}}
\newcommand{\Pic}{\operatorname{Pic}}
\newcommand{\Jac}{\operatorname{Jac}}
\newcommand{\Div}{\operatorname{Div}}
\newcommand{\Prin}{\operatorname{Prin}}
\newcommand{\Prob}{\operatorname{Prob}}
\newcommand{\Gal}{\operatorname{Gal}}
\newcommand{\tors}{\mathrm{tors}}
\newcommand{\id}{\mathrm{id}}

\PassOptionsToPackage{hyphens}{url}
\usepackage[hidelinks,bookmarks=false]{hyperref}

\title[Isogeny and the Picard torsor]{Isogeny, zeta-equivalence, and the Picard torsor\\ in function field Bost--Connes systems}
\author{Ruqing Chen}
\address{GUT Geoservice Inc., Montreal, Canada}
\email{ruqing@hotmail.com}
\thanks{Sources, code and data for the entire series: \url{https://github.com/Ruqing1963/localized-bost-connes}; this paper is in \texttt{papers/XIX-isogeny/}.}
\date{}

\begin{document}

\begin{abstract}
Paper XVIII of the series showed that the partition function of the function field
Bost--Connes system is the Weil zeta function, and asked whether the system can distinguish
two curves with the same zeta function but non-isomorphic Jacobian groups. The answer is yes,
and we locate precisely which layer of the system does it.

The proposed exact sequence $1\to\Jac(C)(\F_q)\to G_{\cP_K}\to\hZ\to1$ is not the right one:
$\Jac(C)(\F_q)=\Pic^0(C)$ is a \emph{quotient} of the compact part
$\Idl^0_K/K^\times$, not a subgroup of $G_{\cP_K}$, and the correct picture is a two-step
filtration with the much larger $\prod_P\cO_P^\times/\overline{\F_q^\times}$ at the bottom.
What is true, and is what the disambiguation needs, is the identity
\[
G_{\cP_K}\big/\overline{U_S}\ \cong\ \widehat{\Pic(C)}\ \cong\ \hZ\times\Jac(C)(\F_q),
\qquad U_S=\prod_P\cO_P^\times,
\]
coming from $\Idl_K/(K^\times U)=\Div(C)/\Prin(C)=\Pic(C)$ together with F.~K. Schmidt's
theorem that $C$ has a divisor of degree $1$. Since $\hZ$ is torsion free,
\[
\Jac(C)(\F_q)=\bigl(G_{\cP_K}/\overline{U_S}\bigr)_{\tors},
\]
so the Jacobian group is recovered from the \emph{pair} $(G_{\cP_K},U_S)$ --- data the
groupoid supplies but the partition function does not.

Hence the disambiguation theorem holds, at the Cartan-pair level: a gauge-equivariant
isomorphism of systems forces $\Jac(C_1)(\F_q)\cong\Jac(C_2)(\F_q)$, while the partition
function forces only $Z_{C_1}=Z_{C_2}$, that is, isogeny. The separation is not vacuous. By
brute-force enumeration over $\F_p$ for $p\le23$ we exhibit many isogeny classes containing
several group structures; for $p=13$ and trace $-2$ there are three,
$\Z/16$, $\Z/2\times\Z/8$ and $\Z/4\times\Z/4$, all with the same zeta function.

Two cautions. Over a function field $S=\cP_K$ is massively \emph{not} norm-separated ---
$\Ns(D)=q^{\deg D}$ identifies all divisors of equal degree --- so Theorem 2.5 of Paper XIII
does not apply at all and the gauge-equivariant hypothesis of Theorem 4.2 of Paper XIV is
not a convenience but a necessity. And the proposed third tier, recovery of $C$ itself by
Torelli, does not follow from anything here: what the system supplies is abelian class field
theory data, and the abelianized Galois group is known not to determine a global field in the
number field case. We record it as open, not proved.
\end{abstract}

\maketitle

\section{The correct filtration}

Let $C/\F_q$ be smooth projective geometrically irreducible, $K=\F_q(C)$, $S=\cP_K$ the set
of closed points. The framework is that of \cite{Main} transported to function fields as in
Paper XVII of the series, including its \S1 convention that $G_{\cP_K}$ denotes the
profinite completion in the degree direction; Papers XIII, XIV, XVII and XVIII, referred to
by numeral, are unpublished companions, and every fact used from them is restated or
verified in place. Function fields have no archimedean places, so
$\Idl_S=\Idl_K$, $K^\times_S=K^\times$ and
\[
G_{\cP_K}=\Idl_K/\overline{K^\times},\qquad U_S=\prod_{P}\cO_P^\times .
\]

\begin{proposition}[Two-step filtration]\label{prop:filtration}
There are exact sequences of topological groups
\begin{align}
1&\longrightarrow \prod_P\cO_P^\times\big/\overline{\F_q^\times}\longrightarrow \Idl^0_K/K^\times\longrightarrow\Pic^0(C)(\F_q)\longrightarrow1,\label{eq:seq1}\\
1&\longrightarrow \Idl^0_K/K^\times\longrightarrow G_{\cP_K}\xrightarrow{\ \deg\ }\hZ\longrightarrow1,\label{eq:seq2}
\end{align}
where $\Idl^0_K$ denotes the degree-zero ideles. The group $\Idl^0_K/K^\times$ is compact.
\end{proposition}

\begin{proof}
Standard class field theory for function fields: \eqref{eq:seq2} is the degree map, whose
image after profinite completion is $\hZ$; \eqref{eq:seq1} is the quotient of the degree-zero
idele class group by the unit ideles, whose cokernel is the degree-zero divisor class group.
\end{proof}

\begin{remark}[The proposed sequence]\label{rem:proposedseq}
The sequence $1\to\Jac(C)(\F_q)\to G_{\cP_K}\to\hZ\to1$ is not correct. By
\eqref{eq:seq1}, $\Jac(C)(\F_q)=\Pic^0(C)$ is a \emph{quotient} of the compact part, not a
subgroup, and the kernel in \eqref{eq:seq2} is $\Idl^0_K/K^\times$, an extension of
$\Pic^0(C)$ by the enormous group $\prod_P\cO_P^\times/\overline{\F_q^\times}$. The Jacobian
group is a subquotient of $G_{\cP_K}$, and exhibiting it requires naming $U_S$.
\end{remark}

\section{The Jacobian group as a torsion subgroup}

\begin{theorem}[Recovery from the pair]\label{thm:torsion}
With $U_S=\prod_P\cO_P^\times\le G_{\cP_K}$,
\begin{equation}\label{eq:picard}
G_{\cP_K}\big/\overline{U_S}\ \cong\ \widehat{\Pic(C)}\ \cong\ \hZ\times\Jac(C)(\F_q),
\end{equation}
and consequently
\[
\Jac(C)(\F_q)\ \cong\ \Bigl(G_{\cP_K}\big/\overline{U_S}\Bigr)_{\tors}.
\]
\end{theorem}

\begin{proof}
$\Idl_K/(K^\times U_S)$ is by definition $\Div(C)/\Prin(C)=\Pic(C)$, and passing to profinite
completions gives the first isomorphism in \eqref{eq:picard}. By F.~K. Schmidt's theorem
\cite{Rosen} $C$
has a divisor of degree $1$, so the degree map $\Pic(C)\to\Z$ splits and
$\Pic(C)\cong\Z\oplus\Pic^0(C)$ with $\Pic^0(C)=\Jac(C)(\F_q)$ finite; completing gives
$\hZ\times\Jac(C)(\F_q)$. Finally $\hZ$ is torsion free, so the torsion subgroup of the
product is the finite factor.
\end{proof}

\begin{corollary}[What each layer sees]\label{cor:layers}
The partition function determines $Z_C(T)$ (Theorem 4.1 of Paper XVIII), hence
$|\Jac(C)(\F_q)|=P(1)$ but only its order. The pair $(G_{\cP_K},U_S)$ determines
$\Jac(C)(\F_q)$ as an abstract finite abelian group, by Theorem~\ref{thm:torsion}.
\end{corollary}

\section{Disambiguation}

\begin{theorem}[Isogeny disambiguation]\label{thm:disambig}
Let $C_1,C_2$ be curves over $\F_q$ with $Z_{C_1}=Z_{C_2}$ but
$\Jac(C_1)(\F_q)\not\cong\Jac(C_2)(\F_q)$. Then:
\begin{enumerate}
\item the two systems have the \emph{same} partition function, so no invariant built from
the KMS temperature dependence alone distinguishes them;
\item nevertheless the systems are not isomorphic \emph{gauge-equivariantly}: there is no
isomorphism $\cA_{C_1,\cP_{K_1}}\to\cA_{C_2,\cP_{K_2}}$ intertwining the full gauge actions
in the sense of Paper XIV (its Definition 4.1) --- the actions of the duals of the grading
groups, of which the circle action of Paper XVIII and the flow $\sigma$ are quotients.
\end{enumerate}
\end{theorem}

\begin{proof}
(1) is Theorem 4.1 of Paper XVIII. For (2), suppose $\Phi$ is such an isomorphism. The
argument of Theorem 4.2 of Paper XIV --- the integer-valued gauge cocycle and its bisections
--- uses only the \'etale groupoid picture and the freeness of the grading monoid, and
applies verbatim over a function field; it gives $\Phi(C(Y_{S_1}))=C(Y_{S_2})$ together
with a monoid isomorphism $\theta:J_{S_1}\to J_{S_2}$, i.e.\ a $\theta$-equivariant
homeomorphism $h:Y_{S_1}\to Y_{S_2}$ of the semigroup dynamical systems. Now consider the
\emph{deep stratum}
\[
Y_\infty:=\bigcap_{0\ne D\in J_S}D\cdot Y_S=\{[x,g]:x=0\}\ \cong\ G_{\cP_K}\big/\overline{U_S},
\]
canonically defined from the dynamics, on which $D\in J_S$ acts by the translation
$[0,g]\mapsto[0,\sigma_Dg]$: the deep stratum is a torsor over the compactified Picard group
$G_{\cP_K}/\overline{U_S}\cong\widehat{\Pic(C)}$ of Theorem~\ref{thm:torsion}, and the
translations $\{\sigma_D\}$ run through the effective classes, which generate $\Pic(C)$ and
are dense in the completion. The closure of the translation monoid in
$\mathrm{Homeo}(Y_\infty)$ is therefore a compact group acting simply transitively --- a
closed subsemigroup of a compact group is a subgroup --- and it is isomorphic, as a
topological group, to $G_{\cP_K}/\overline{U_S}$. Since $h$ maps $Y_{\infty,1}$ onto
$Y_{\infty,2}$ intertwining the translation monoids, it intertwines their closures, giving
a topological group isomorphism
$G_{\cP_{K_1}}/\overline{U_{S_1}}\cong G_{\cP_{K_2}}/\overline{U_{S_2}}$. Taking torsion
subgroups, Theorem~\ref{thm:torsion} forces $\Jac(C_1)(\F_q)\cong\Jac(C_2)(\F_q)$, contrary
to hypothesis. (Note that the homeomorphism type of $Y_\infty$ alone would see nothing: all
these spaces are Cantor sets, and indeed $|\Jac|$ is already matched by (1); it is the
translation structure that carries the group.)
\end{proof}

\begin{remark}[Gauge-equivariance is necessary here]\label{rem:necessary}
Over a function field $\Ns(D)=q^{\deg D}$, so $\Ns$ identifies all divisors of the same
degree and $S=\cP_K$ is as far from norm-separated as possible: the relation lattice
$L(S)$ of Paper XIV (its Definition 1.2) has infinite rank. Consequently
Theorem 1.4 of Paper XIII, which identifies $C(Y_S)$ with $\cA^\sigma$ under
norm-separation, is unavailable, and
$\cA^\sigma$ is a large non-commutative algebra --- as is the fixed-point algebra of the
circle action, which coincides with it. Theorem~\ref{thm:disambig} therefore rests on
Theorem 4.2 of Paper XIV and its \emph{full} gauge-equivariant hypothesis; equivariance for
the circle action alone is the same as $\sigma$-equivariance and does not feed that theorem.
In this setting the hypothesis is not a
technical convenience but the only route available.
\end{remark}

\begin{example}[Explicit isogeny classes with several group structures]\label{ex:explicit}
Enumerating all $y^2=x^3+ax+b$ over $\F_p$ and grouping by trace $t$ --- equal $t$ being
equivalent to equal zeta function, hence to isogeny \cite{Tate} --- gives, for instance:
\[
\begin{array}{lll}
p & t & \text{group structures of }E(\F_p)\\\hline
11 & 0 & \Z/12,\ \Z/2\times\Z/6\\
13 & -2 & \Z/16,\ \Z/2\times\Z/8,\ \Z/4\times\Z/4\\
13 & 5 & \Z/9,\ \Z/3\times\Z/3\\
17 & 2 & \Z/16,\ \Z/2\times\Z/8,\ \Z/4\times\Z/4\\
19 & -7 & \Z/27,\ \Z/3\times\Z/9\\
23 & -8 & \Z/32,\ \Z/2\times\Z/16
\end{array}
\]
Each row is a single isogeny class, so a single partition function, containing curves whose
systems are pairwise non-isomorphic by Theorem~\ref{thm:disambig}. The phenomenon is
abundant, not exceptional: it occurred for most traces at every prime $11\le p\le23$ examined.
\end{example}

\section{The hierarchy, with the third tier open}

\begin{theorem}[Two tiers]\label{thm:tiers}
For the function field system $(\cA_{C,\cP_K},\T,\sigma)$:
\begin{itemize}
\item \textbf{Tier 1, thermodynamic.} The partition function recovers $Z_C(T)$, hence the
Weil polynomial, the genus, all point counts $\#C(\F_{q^n})$, and $|\Jac(C)(\F_q)|$. Two
systems have the same partition function iff the curves are zeta-equivalent, that is iff
their Jacobians are isogenous over $\F_q$ \cite{Tate}.
\item \textbf{Tier 2, Cartan pair.} A gauge-equivariant isomorphism recovers the diagonal
with its semigroup action, hence the deep stratum with its translation structure, hence the
compactified Picard group and $\Jac(C)(\F_q)$ as an abstract group, strictly refining Tier 1
by Example~\ref{ex:explicit}.
\end{itemize}
\end{theorem}

\begin{remark}[The third tier is not established]\label{rem:tier3}
The proposal's third tier --- recovery of $C$ itself, ``the Torelli geometric isomorphism'' ---
does not follow from what is proved here, and we do not believe it follows from abelian data
at all. What the system supplies is the abelianized Galois group with its Frobenius classes:
class field theory data. In the number field case the abelianized absolute Galois group is
known not to determine the field --- Onabe, and Angelakis--Stevenhagen \cite{AS}, exhibit infinitely
many imaginary quadratic fields sharing the same $\Gal(K^{ab}/K)$ --- so no purely abelian
invariant can give a Torelli statement. Our data is richer than the bare group, carrying the
Frobenius classes indexed by $S$ together with their degrees; whether that suffices is
exactly the question addressed in the literature on characterizing global fields by their
$L$-series, which we have not been able to consult here. We state Tier 3 as open and make no
claim.
\end{remark}

\begin{remark}[What Tier 2 does and does not give]\label{rem:tier2limits}
Theorem~\ref{thm:torsion} recovers $\Jac(C)(\F_q)$ as an abstract finite abelian group. It
does not recover the principal polarization, which is what Torelli's theorem requires (for
the arithmetic of pairings and polarizations over finite fields see \cite{Howe}), nor
the Galois module structure of the $\ell$-adic Tate module beyond its characteristic
polynomial. So even a complete solution at Tier 2 would leave a gap before Tier 3, and the
gap is exactly the polarization.
\end{remark}

\section{Assessment}

\[
\begin{array}{lll}
\text{Statement} & \text{Status} & \text{Reference}\\\hline
1\to\Jac(\F_q)\to G_{\cP_K}\to\hZ\to1 & \textbf{false: wrong kernel} & \text{Rem.~1.2}\\
\text{two-step filtration} & \text{true} & \text{Prop.~1.1}\\
G_{\cP_K}/\overline{U_S}\cong\hZ\times\Jac(\F_q) & \textbf{true} & \text{Thm.~2.1}\\
\Jac(\F_q)=\bigl(G_{\cP_K}/\overline{U_S}\bigr)_{\tors} & \textbf{true} & \text{Thm.~2.1}\\
\text{partition function}\iff\text{isogeny} & \text{true} & \text{Thm.~4.1}\\
\text{system distinguishes isogenous curves} & \textbf{yes, at the Cartan pair} & \text{Thm.~3.1}\\
S=\cP_K\ \text{norm-separated} & \text{no; }L(S)\ \text{of infinite rank} & \text{Rem.~3.2}\\
\text{gauge-equivariance needed} & \text{yes, necessarily} & \text{Rem.~3.2}\\
\text{Tier 3: }C_1\cong C_2 & \textbf{open; not from abelian data} & \text{Rem.~4.2}
\end{array}
\]

The question left open in \S5 of Paper XVIII is answered: zeta-equivalence is
\emph{not} the equivalence realized by these systems. The partition function sees exactly the
isogeny class, and the Cartan pair sees strictly more, namely the isomorphism class of
$\Jac(C)(\F_q)$ as a group. The mechanism is the identity
$\Idl_K/(K^\times U)=\Pic(C)$, which is class field theory rather than operator algebra; the
operator-algebraic content is only that a gauge-equivariant isomorphism recovers the
diagonal with its semigroup action, which is Theorem 4.2 of Paper XIV --- the passage from
there to the Picard group is the translation structure of the deep stratum.

We are left with a clean picture of the layers and an honest gap at the top. Thermodynamics
gives the isogeny class; the Cartan pair gives the Jacobian group; and the curve itself is
out of reach of both, because the polarization is nowhere in the data and because abelian
class field theory is known to be insufficient for reconstructing global fields.

Three questions. Does the pair $(G_S,U_S)$ with the Frobenius classes determine more than
$\Jac(C)(\F_q)$ --- for instance the full group $\Pic(C)$ with its effective cone, which is
closer to geometry? Second, Theorem~\ref{thm:torsion} used F.~K. Schmidt to split the degree
map; is the splitting canonical enough that the isomorphism $\Jac(C)(\F_q)\cong(\cdot)_\tors$
is functorial, or does the choice of degree-one divisor matter for the disambiguation?
Third, and this is the natural next target, is there a curve pair with the same zeta
\emph{and} the same $\Jac(\F_q)$ but non-isomorphic curves --- for which Tier 2 would also
fail, and the question of Tier 3 would become sharp rather than speculative?

\begin{thebibliography}{9}
\bibitem{AS} A. Angelakis, P. Stevenhagen, \emph{Imaginary quadratic fields with isomorphic abelian Galois groups}, ANTS X, Open Book Ser. 1 (2013), 21--39.
\bibitem{Howe} E. W. Howe, \emph{The Weil pairing and the Hilbert symbol}, Math. Ann. 305 (1996), 387--392.
\bibitem{Rosen} M. Rosen, \emph{Number Theory in Function Fields}, Graduate Texts in Math. 210, Springer, 2002.
\bibitem{Tate} J. Tate, \emph{Endomorphisms of abelian varieties over finite fields}, Invent. Math. 2 (1966), 134--144.
\bibitem{Main} R. Chen, \emph{KMS state uniqueness and phase transitions in localized Bost--Connes systems}, preprint, 2026, \newline\url{https://doi.org/10.5281/zenodo.22152101}.
\end{thebibliography}

\end{document}
